% ============================================================
%  CHAPTER 6 – TEST PLAN
% ============================================================
\chapter{Test Plan}
\label{ch:test-plan}
This chapter discusses the individual subsystem verification test plans (Section~\ref{sec:individual-tests}), subsystem integration testing (Section~\ref{sec:integration-tests}), and full-system validation (Section~\ref{sec:system-validation}) to be implemented to ensure the functioning of each subsystem as a whole. The persons responsible for collaboration on the test plan are the 2 persons who have the role in the subsection title. Testing is divided into individual subsystem verification,
subsystem integration testing, and full-system validation.

Each test shall specify the initial conditions, test procedure,
expected result, and acceptance criteria. The actual result and
test status shall be recorded during execution.

\section{Test Procedure and Acceptance Criteria}
The following information shall be documented for each test:
\begin{itemize}
    \item \textbf{Test ID:} A unique identifier for the test.
    \item \textbf{Objective:} The functionality or requirement
    being verified.
    \item \textbf{Initial conditions:} The required system
    configuration before testing.
    \item \textbf{Procedure:} The steps followed to perform the test.
    \item \textbf{Expected result:} The anticipated system response.
    \item \textbf{Acceptance criteria:} The measurable conditions
    required for the test to pass.
    \item \textbf{Actual result:} The observed outcome.
    \item \textbf{Status:} Pass, Fail, or Not Tested.
\end{itemize}

\section{Individual Test Plans}
\label{sec:individual-tests}

Each subsystem shall first be tested independently, where
practical, to establish that its functionality meets the
specified design requirements before integration testing.


\subsection{EERI I}
\label{subsec:eeri-i-tests}

\textbf{Subsystem function:} [Insert the function of EERI I.]

\textbf{Test setup:} [Insert the required equipment, configuration,
and initial conditions.]

\begin{table}[htbp]
    \centering
    \caption{EERI I test plan}
    \label{tab:eeri-i-tests}
    \renewcommand{\arraystretch}{1.4}
    \begin{tabularx}{\textwidth}{|l|X|X|X|X|}
        \hline
        \textbf{Test ID} &
        \textbf{Test description} &
        \textbf{Procedure} &
        \textbf{Expected result} &
        \textbf{Status} \\
        \hline
        EI-01 & [Functional test] &
        [Test procedure] &
        [Expected behaviour] &
        [Pass/Fail/NT] \\
        \hline
        EI-02 & [Performance test] &
        [Test procedure] &
        [Acceptance criterion] &
        [Pass/Fail/NT] \\
        \hline
        EI-03 & [Safety or fault test] &
        [Test procedure] &
        [Expected safe response] &
        [Pass/Fail/NT] \\
        \hline
    \end{tabularx}
\end{table}


\subsection{EERI II}
\label{subsec:eeri-ii-tests}

\textbf{Subsystem function:} [Insert the function of EERI II.]

\textbf{Test setup:} [Insert the required equipment, configuration,
and initial conditions.]

\begin{table}[htbp]
    \centering
    \caption{EERI II test plan}
    \label{tab:eeri-ii-tests}
    \renewcommand{\arraystretch}{1.4}
    \begin{tabularx}{\textwidth}{|l|X|X|X|X|}
        \hline
        \textbf{Test ID} &
        \textbf{Test description} &
        \textbf{Procedure} &
        \textbf{Expected result} &
        \textbf{Status} \\
        \hline
        EII-01 & [Functional test] &
        [Test procedure] &
        [Expected behaviour] &
        [Pass/Fail/NT] \\
        \hline
        EII-02 & [Performance test] &
        [Test procedure] &
        [Acceptance criterion] &
        [Pass/Fail/NT] \\
        \hline
        EII-03 & [Safety or fault test] &
        [Test procedure] &
        [Expected safe response] &
        [Pass/Fail/NT] \\
        \hline
    \end{tabularx}
\end{table}


\subsection{MCTR I}
\label{subsec:mctr-i-tests}

\textbf{Subsystem function:} [Insert the function of MCTR I.]

\textbf{Test setup:} [Insert the required equipment, configuration,
and initial conditions.]

\begin{table}[htbp]
    \centering
    \caption{MCTR I test plan}
    \label{tab:mctr-i-tests}
    \renewcommand{\arraystretch}{1.4}
    \begin{tabularx}{\textwidth}{|l|X|X|X|X|}
        \hline
        \textbf{Test ID} &
        \textbf{Test description} &
        \textbf{Procedure} &
        \textbf{Expected result} &
        \textbf{Status} \\
        \hline
        MCI-01 & [Functional test] &
        [Test procedure] &
        [Expected behaviour] &
        [Pass/Fail/NT] \\
        \hline
        MCI-02 & [Performance test] &
        [Test procedure] &
        [Acceptance criterion] &
        [Pass/Fail/NT] \\
        \hline
        MCI-03 & [Safety or fault test] &
        [Test procedure] &
        [Expected safe response] &
        [Pass/Fail/NT] \\
        \hline
    \end{tabularx}
\end{table}


\subsection{MCTR II}
\label{subsec:mctr-ii-tests}

\textbf{Subsystem function:} [Insert the function of MCTR II.]

\textbf{Test setup:} [Insert the required equipment, configuration,
and initial conditions.]

\begin{table}[htbp]
    \centering
    \caption{MCTR II test plan}
    \label{tab:mctr-ii-tests}
    \renewcommand{\arraystretch}{1.4}
    \begin{tabularx}{\textwidth}{|l|X|X|X|X|}
        \hline
        \textbf{Test ID} &
        \textbf{Test description} &
        \textbf{Procedure} &
        \textbf{Expected result} &
        \textbf{Status} \\
        \hline
        MCII-01 & [Functional test] &
        [Test procedure] &
        [Expected behaviour] &
        [Pass/Fail/NT] \\
        \hline
        MCII-02 & [Performance test] &
        [Test procedure] &
        [Acceptance criterion] &
        [Pass/Fail/NT] \\
        \hline
        MCII-03 & [Safety or fault test] &
        [Test procedure] &
        [Expected safe response] &
        [Pass/Fail/NT] \\
        \hline
    \end{tabularx}
\end{table}


\subsection{REII}
\label{subsec:reii-tests}

\textbf{Subsystem function:} [Insert the function of REII.]

\textbf{Test setup:} [Insert the required equipment, configuration,
and initial conditions.]

\begin{table}[htbp]
    \centering
    \caption{REII test plan}
    \label{tab:reii-tests}
    \renewcommand{\arraystretch}{1.4}
    \begin{tabularx}{\textwidth}{|l|X|X|X|X|}
        \hline
        \textbf{Test ID} &
        \textbf{Test description} &
        \textbf{Procedure} &
        \textbf{Expected result} &
        \textbf{Status} \\
        \hline
        RE-01 & [Functional test] &
        [Test procedure] &
        [Expected behaviour] &
        [Pass/Fail/NT] \\
        \hline
        RE-02 & [Performance test] &
        [Test procedure] &
        [Acceptance criterion] &
        [Pass/Fail/NT] \\
        \hline
        RE-03 & [Safety or fault test] &
        [Test procedure] &
        [Expected safe response] &
        [Pass/Fail/NT] \\
        \hline
    \end{tabularx}
\end{table}

\section{Subsystem Interface Analysis}
\label{sec:interface-analysis}

The subsystem interface analysis identifies the direct interactions
between subsystems. These interfaces may involve mechanical loads,
electrical power, control signals, communication, or physical
connections, as applicable to the design.

\begin{table}[htbp]
    \centering
    \caption{Subsystem interface matrix}
    \label{tab:interface-matrix}
    \renewcommand{\arraystretch}{1.4}
    \begin{tabularx}{\textwidth}{|X|c|c|c|c|c|}
        \hline
        \textbf{Subsystem} &
        \textbf{EERI I} &
        \textbf{EERI II} &
        \textbf{MCTR I} &
        \textbf{MCTR II} &
        \textbf{REII} \\
        \hline
        EERI I   & -- & [X] & [X] & [X] & [X] \\
        \hline
        EERI II  & [X] & -- & [X] & [X] & [X] \\
        \hline
        MCTR I   & [X] & [X] & -- & [X] & [X] \\
        \hline
        MCTR II  & [X] & [X] & [X] & -- & [X] \\
        \hline
        REII     & [X] & [X] & [X] & [X] & -- \\
        \hline
    \end{tabularx}
\end{table}

In Table~\ref{tab:interface-matrix}, replace \texttt{[X]} with
\texttt{X} where a direct interface exists and leave the cell
blank where no direct interface exists. The matrix shall be
used to determine the required integration tests.

\section{Subsystem Integration Test Plan}
\label{sec:integration-tests}

Integration testing verifies that connected subsystems operate
together correctly and that their interfaces satisfy the design
requirements. Testing shall be performed after the relevant
individual subsystem tests have been completed.

\begin{table}[htbp]
    \centering
    \caption{Subsystem integration test plan}
    \label{tab:integration-tests}
    \renewcommand{\arraystretch}{1.4}
    \begin{tabularx}{\textwidth}{|l|X|X|X|X|}
        \hline
        \textbf{Test ID} &
        \textbf{Subsystems} &
        \textbf{Test objective} &
        \textbf{Expected result} &
        \textbf{Status} \\
        \hline
        INT-01 & [Subsystem A and B] &
        Verify the interface between the subsystems &
        [Specified interface behaviour] &
        [Pass/Fail/NT] \\
        \hline
        INT-02 & [Subsystem A and C] &
        Verify coordinated operation &
        [Specified response] &
        [Pass/Fail/NT] \\
        \hline
        INT-03 & [Subsystem B and D] &
        Verify signal, load, or power transfer, as applicable &
        [Specified transfer behaviour] &
        [Pass/Fail/NT] \\
        \hline
        INT-04 & [Relevant subsystems] &
        Verify a multi-subsystem operating sequence &
        [Specified sequence completed correctly] &
        [Pass/Fail/NT] \\
        \hline
        INT-05 & [Relevant subsystems] &
        Verify fault propagation and safe response &
        [Specified fault-handling behaviour] &
        [Pass/Fail/NT] \\
        \hline
    \end{tabularx}
\end{table}

For each integration test, record the initial conditions, test
procedure, measured or observed results, and applicable
acceptance criteria. Tests shall cover every required interface
identified in Table~\ref{tab:interface-matrix}.


\section{Full-System Validation}
\label{sec:system-validation}

Full-system validation shall be conducted after the individual
subsystem and integration tests have been completed. The purpose
is to determine whether the complete gantry crane system satisfies
its overall functional, performance, and safety requirements.

\begin{table}[htbp]
    \centering
    \caption{Full-system validation plan}
    \label{tab:system-validation}
    \renewcommand{\arraystretch}{1.4}
    \begin{tabularx}{\textwidth}{|l|X|X|X|}
        \hline
        \textbf{Test ID} &
        \textbf{Test scenario} &
        \textbf{Acceptance criteria} &
        \textbf{Status} \\
        \hline
        SYS-01 & Normal startup and shutdown &
        [System starts and stops as specified] &
        [Pass/Fail/NT] \\
        \hline
        SYS-02 & Complete operating cycle &
        [Required sequence is completed] &
        [Pass/Fail/NT] \\
        \hline
        SYS-03 & Operation under specified load &
        [Performance remains within limits] &
        [Pass/Fail/NT] \\
        \hline
        SYS-04 & Fault handling and emergency stop &
        [Required safe response is achieved] &
        [Pass/Fail/NT] \\
        \hline
        SYS-05 & Repeated operation &
        [Required repeatability or reliability criteria] &
        [Pass/Fail/NT] \\
        \hline
    \end{tabularx}
\end{table}


\section{Test Results Summary}
\label{sec:test-results-summary}

The results of the subsystem, integration, and full-system tests
shall be consolidated to assess overall system readiness.

\begin{table}[htbp]
    \centering
    \caption{Summary of test outcomes}
    \label{tab:test-summary}
    \renewcommand{\arraystretch}{1.4}
    \begin{tabularx}{\textwidth}{|X|c|c|c|c|}
        \hline
        \textbf{Test category} &
        \textbf{Passed} &
        \textbf{Failed} &
        \textbf{Not tested} &
        \textbf{Total} \\
        \hline
        Individual subsystem tests & [\,] & [\,] & [\,] & [\,] \\
        \hline
        Integration tests & [\,] & [\,] & [\,] & [\,] \\
        \hline
        Full-system validation & [\,] & [\,] & [\,] & [\,] \\
        \hline
        \textbf{Total} & [\,] & [\,] & [\,] & [\,] \\
        \hline
    \end{tabularx}
\end{table}

Any failed or incomplete tests shall be documented together with
the observed behaviour, likely cause, corrective action, and
retesting requirements. The system shall only be declared to have
passed validation when the applicable acceptance criteria have
been met.
